\documentclass[12pt]{article}
\usepackage[margin=1in]{geometry}
\usepackage[T1]{fontenc}
\usepackage{bold-extra}
\usepackage{xcolor}
\definecolor{garnet}{HTML}{73000A}

\setlength{\parindent}{0pt}
\setlength{\parskip}{0.5em}

\begin{document}

{\color{garnet}\large\textbf{\textsc{SC Space Grant Personal Essay}}}\\[-0.8em]
{\color{garnet}\rule{\textwidth}{0.5pt}}\\[0.2em]
It should not have flown. A gyrocopter built from car transmissions, helicopter rotors, airplane controls, and snowmobile skis. The pilot at USC Fly Club had assembled it himself from parts never meant to share an airframe. But it flew. He did not need aerospace-grade components. He needed to understand the essential physics of flight and find parts that satisfied those requirements. That distinction has shaped how I approach every problem since.

My Honors Introduction to Aerospace Engineering course formalized this intuition. As the class integrated concepts from statics, controls, system dynamics, and aerodynamics, I saw how aerospace engineering requires balancing theoretical models with practical data and constraints. This experience shifted my interest from a general fascination with aerospace systems to a focused curiosity about developing effective representations of complex systems that enable meaningful analysis and innovation.

I built a simplified simulation of aircraft powertrains using MATLAB and Simscape to understand how complex systems can be represented efficiently without sacrificing physical relevance. In parallel, I explored aircraft modeling using CAD software, which further strengthened my understanding of system structure and component interaction. These projects emphasized the importance of selecting appropriate assumptions and demonstrated that effective simplification requires knowing which details matter and which can be abstracted away. This work motivated me to apply similar modeling principles to subsequent projects, including the ongoing design of a hovercraft.

In Dr.\ Yi Wang's lab, I apply these principles to radar systems for autonomous maritime navigation. I have fine-tuned object detection models for deployment on embedded hardware, built pipelines that convert raw radar returns into usable imagery, and designed sensor mounts for research vessel integration. I present progress to Department of Defense sponsors monthly. A persistent challenge is validating navigation algorithms without expensive field campaigns. My proposed research addresses this gap: developing an open-source radar simulator that models terrain occlusion through geometric ray-casting rather than costly electromagnetic computation. By validating against real offshore measurements, the simulator will establish quantitative accuracy bounds -- showing researchers when simplified models suffice and when higher fidelity is required. This makes algorithm validation accessible to teams without the budget for full physics-based simulation.

Recent advances by Boom Supersonic and NASA's X-59 QUESST program illustrate how aerospace progress increasingly depends on reexamining long-standing assumptions and improving how complex systems are modeled and validated. Before any autonomous system flies, its algorithms are tested thousands of times in simulation -- NASA's Mars 2020 lander validated terrain-relative navigation this way, achieving 5-meter accuracy after extensive ground testing.

Support from the NASA South Carolina Space Grant would allow me to further refine these modeling and systems-level skills through research-driven projects. By providing the time, mentorship, and resources necessary for focused exploration, Space Grant would prepare me to contribute meaningfully to advanced aerospace and space exploration efforts, where effective system representation and validation are essential to innovation.

\end{document}
